% !TeX program = xelatex
% =============================================================
% 模板 09：Elsevier 期刊投稿（elsarticle，单栏 preprint/review）
% 编译： scripts\build.ps1 en-journal-els
% 版式： A4 单栏 12pt，行号（review 选项），Elsevier 官方 elsarticle.cls
% 用途： SCI 期刊初稿、投稿前排版；正式投稿按目标期刊 Guide for Authors 调整
% =============================================================
\documentclass[preprint,review,12pt]{elsarticle}

\usepackage{amsmath}
\usepackage{fontspec}
\setmainfont{TeX Gyre Termes}
\usepackage{unicode-math}
\setmathfont{TeX Gyre Termes Math}
\usepackage{graphicx}
\usepackage{booktabs}
\usepackage{multirow}
\usepackage{xcolor}
\usepackage[hidelinks]{hyperref}
% ------------------------- 审稿用行号 -------------------------
% elsarticle 的 review 选项只改行距、不加行号；行号要自己开。
% 正式排版（camera-ready）时删掉下面两行即可。
\usepackage{lineno}
\raggedbottom   % 单栏 preprint 常用：页面底部不强制撑满，换页更自然

\journal{Journal of Placeholder Studies}

\begin{document}

\begin{frontmatter}

% ------------------------- 标题与作者（改这里） -------------------------
\title{Placeholder Title: A Sample Elsevier Journal Manuscript for Template Testing}

\author[inst1]{San Zhang\corref{cor1}}
\ead{zhangsan@example.edu.cn}
\cortext[cor1]{Corresponding author.}

\author[inst2]{Si Li}
\ead{lisi@example.edu.cn}

\affiliation[inst1]{organization={School of Electrical Engineering, XX University},
  addressline={No. 1 Placeholder Road},
  city={Beijing},
  postcode={100000},
  country={China}}

\affiliation[inst2]{organization={School of Automation, XX University},
  addressline={No. 2 Placeholder Road},
  city={Beijing},
  postcode={100000},
  country={China}}

% ------------------------- 摘要与关键词 -------------------------
\begin{abstract}
This manuscript is a placeholder for an Elsevier journal submission. It demonstrates the
front matter, a numbered equation, a table, and a bibliography entry so that the layout
can be checked before the real content is written. Replace every sentence here with your
own 150--250 word abstract. A journal abstract should state the problem, the method, the
main quantitative results, and the conclusion, without citations or undefined abbreviations.
\end{abstract}

\begin{keyword}
placeholder \sep template \sep Elsevier \sep elsarticle \sep journal manuscript
\end{keyword}

\end{frontmatter}

% 正文才加行号：放在 frontmatter 之后，避免标题区脚注触发行号溢出
\linenumbers

% ------------------------- 正文 -------------------------
\section{Introduction}
\label{sec:intro}

This is the introduction placeholder. A journal introduction usually covers the research
background, a critical review of related work~\cite{lecun2015deep,he2016resnet}, the
research gap, and a clearly stated list of contributions at the end of the section.

\section{Materials and Methods}
\label{sec:method}

Let $x_i$ denote the $i$-th input feature and $w_i$ its weight. The weighted sum is
\begin{equation}
  y=\sum_{i=1}^{n} w_i x_i + b,
  \label{eq:linear}
\end{equation}
where $b$ is the bias term. Equation~\eqref{eq:linear} is the equation placeholder.

The parameter settings are listed in Table~\ref{tab:params}.

\begin{table}[htbp]
  \centering
  \caption{Placeholder experimental parameters.}
  \label{tab:params}
  \begin{tabular}{lccc}
    \toprule
    Parameter & Symbol & Value & Unit \\
    \midrule
    Learning rate & $\eta$    & 0.01 & --  \\
    Epochs        & $E$       & 100  & --  \\
    Batch size    & $B$       & 32   & --  \\
    Weight decay  & $\lambda$ & 0.10 & --  \\
    \bottomrule
  \end{tabular}
\end{table}

\section{Results and Discussion}
\label{sec:results}

With the settings in Table~\ref{tab:params}, the placeholder model reaches 92.4\% accuracy
on the validation split. Report mean and standard deviation over repeated runs, state the
statistical test used, and discuss both the improvement and its limitations. Figures should
be vector graphics (PDF/EPS) at final print size.

\section{Conclusions}
\label{sec:conclusion}

This template compiles end-to-end with XeLaTeX and BibTeX. Replace the placeholder text,
figures, and references with your own content before submission, and follow the target
journal's Guide for Authors for reference style, highlights, and graphical abstract.

% ------------------------- 声明（按期刊要求增删） -------------------------
\newpage   % 声明与参考文献单独起页，避免 CRediT 段落跨页
\section*{CRediT authorship contribution statement}
\textbf{San Zhang:} Conceptualization, Methodology, Writing -- original draft.
\textbf{Si Li:} Validation, Writing -- review \& editing.

\section*{Declaration of competing interest}
The authors declare that they have no known competing financial interests or personal
relationships that could have appeared to influence the work reported in this paper.

\section*{Data availability}
Data will be made available on request.

% ------------------------- 参考文献 -------------------------
\bibliographystyle{elsarticle-num}
\bibliography{refs}

\end{document}
